\documentclass{amsart}
% For dealing with references we use the comment environment
\usepackage{verbatim}
\newenvironment{reference}{\comment}{\endcomment}
%\newenvironment{reference}{}{}
\newenvironment{slogan}{\comment}{\endcomment}
\newenvironment{history}{\comment}{\endcomment}

% For commutative diagrams we use Xy-pic
\usepackage[all]{xy}

% We use 2cell for 2-commutative diagrams.
\xyoption{2cell}
\UseAllTwocells

% We use multicol for the list of chapters between chapters
\usepackage{multicol}

% This is generally recommended for better output
\usepackage{lmodern}
\usepackage[T1]{fontenc}

% For cross-file-references

% Package for hypertext links:
\usepackage{hyperref}

% For any local file, say "hello.tex" you want to link to please

% Theorem environments.
%
\theoremstyle{plain}
\newtheorem{theorem}[subsection]{Theorem}
\newtheorem{proposition}[subsection]{Proposition}
\newtheorem{lemma}[subsection]{Lemma}

\theoremstyle{definition}
\newtheorem{definition}[subsection]{Definition}
\newtheorem{example}[subsection]{Example}
\newtheorem{exercise}[subsection]{Exercise}
\newtheorem{situation}[subsection]{Situation}

\theoremstyle{remark}
\newtheorem{remark}[subsection]{Remark}
\newtheorem{remarks}[subsection]{Remarks}

\numberwithin{equation}{subsection}

% Macros
%
\def\lim{\mathop{\mathrm{lim}}\nolimits}
\def\colim{\mathop{\mathrm{colim}}\nolimits}
\def\Spec{\mathop{\mathrm{Spec}}}
\def\Hom{\mathop{\mathrm{Hom}}\nolimits}
\def\Ext{\mathop{\mathrm{Ext}}\nolimits}
\def\SheafHom{\mathop{\mathcal{H}\!\mathit{om}}\nolimits}
\def\SheafExt{\mathop{\mathcal{E}\!\mathit{xt}}\nolimits}
\def\Sch{\mathit{Sch}}
\def\Mor{\mathop{\mathrm{Mor}}\nolimits}
\def\Ob{\mathop{\mathrm{Ob}}\nolimits}
\def\Sh{\mathop{\mathit{Sh}}\nolimits}
\def\NL{\mathop{N\!L}\nolimits}
\def\CH{\mathop{\mathrm{CH}}\nolimits}
\def\proetale{{pro\text{-}\acute{e}tale}}
\def\etale{{\acute{e}tale}}
\def\QCoh{\mathit{QCoh}}
\def\Ker{\mathop{\mathrm{Ker}}}
\def\Im{\mathop{\mathrm{Im}}}
\def\Coker{\mathop{\mathrm{Coker}}}
\def\Coim{\mathop{\mathrm{Coim}}}

% Boxtimes
%
\DeclareMathSymbol{\boxtimes}{\mathbin}{AMSa}{"02}

%
% Macros for moduli stacks/spaces
%
\def\QCohstack{\mathcal{QC}\!\mathit{oh}}
\def\Cohstack{\mathcal{C}\!\mathit{oh}}
\def\Spacesstack{\mathcal{S}\!\mathit{paces}}
\def\Quotfunctor{\mathrm{Quot}}
\def\Hilbfunctor{\mathrm{Hilb}}
\def\Curvesstack{\mathcal{C}\!\mathit{urves}}
\def\Polarizedstack{\mathcal{P}\!\mathit{olarized}}
\def\Complexesstack{\mathcal{C}\!\mathit{omplexes}}
% \Pic is the operator that assigns to X its picard group, usage \Pic(X)
% \Picardstack_{X/B} denotes the Picard stack of X over B
% \Picardfunctor_{X/B} denotes the Picard functor of X over B
\def\Pic{\mathop{\mathrm{Pic}}\nolimits}
\def\Picardstack{\mathcal{P}\!\mathit{ic}}
\def\Picardfunctor{\mathrm{Pic}}
\def\Deformationcategory{\mathcal{D}\!\mathit{ef}}

\setlength{\emergencystretch}{3em}
\begin{document}
\title[Stacks Project, Tag 05F1]{470 --- Relative associated points persist generically along their closures for finite-type morphisms}
\maketitle
\noindent Copyright (C) 2005--2025 Johan de Jong. Stacks Project authors. Source: \href{https://stacks.math.columbia.edu/tag/05F1}{Tag 05F1}.

\noindent Editorial difficulty note (not author text). Difficulty H1: One must turn a fibrewise associated prime into a globally embedded cyclic module after generic flatness. A second generic-flatness step makes the cokernel flat so that the injection survives every fibre, allowing all generic points of the resulting fibre supports to be recognized as associated points..

\noindent Editorial provenance note (not author text). History: excerpt from revision \texttt{a0444\allowbreak 6e57e\allowbreak c1fbc\allowbreak 252a8\allowbreak 71afc\allowbreak ec775\allowbreak 2fb28\allowbreak 07b14}, more-morphisms.tex, lines 6840--6905. Reference presentation was adapted; original-author context and core support blocks were retained or added. The mathematical wording is unchanged. Licensed under GNU FDL 1.2 or later; no invariant sections or cover texts. See ../licenses/COPYING. Context blocks reproduce author definitions and supporting results; they are not additional counted problems.

\begin{definition}
\label{divisors-definition-relative-assassin}
Let $f : X \to S$ be a morphism of schemes.
Let $\mathcal{F}$ be a quasi-coherent $\mathcal{O}_X$-module.
The {\it relative assassin of $\mathcal{F}$ in $X$ over $S$}
is the set
$$
\text{Ass}_{X/S}(\mathcal{F}) =
\bigcup\nolimits_{s \in S} \text{Ass}_{X_s}(\mathcal{F}_s)
$$
where $\mathcal{F}_s = (X_s \to X)^*\mathcal{F}$ is the restriction
of $\mathcal{F}$ to the fibre of $f$ at $s$.
\end{definition}

\begin{lemma}
\label{lemma-relative-assassin-in-neighbourhood}
Let $f : X \to S$ be a morphism of schemes.
Let $\mathcal{F}$ be a quasi-coherent $\mathcal{O}_X$-module.
Let $\xi \in \text{Ass}_{X/S}(\mathcal{F})$ and set
$Z = \overline{\{\xi\}} \subset X$.
If $f$ is locally of finite type and $\mathcal{F}$ is a
finite type $\mathcal{O}_X$-module, then there exists a nonempty
open $V \subset Z$ such that for every $s \in f(V)$ the generic
points of $V_s$ are elements of $\text{Ass}_{X/S}(\mathcal{F})$.
\end{lemma}

\begin{proof}
We may replace $S$ by an affine open neighbourhood of $f(\xi)$
and $X$ by an affine open neighbourhood of $\xi$. Hence we may assume
$S = \Spec(A)$, $X = \Spec(B)$ and that $f$ is given by
the finite type ring map $A \to B$, see
Morphisms, Lemma \href{https://stacks.math.columbia.edu/tag/01T2}{lemma locally finite type characterize (Tag 01T2)}.
Moreover, we may write $\mathcal{F} = \widetilde{M}$ for some finite
$B$-module $M$, see
Properties, Lemma \href{https://stacks.math.columbia.edu/tag/01PB}{lemma finite type module (Tag 01PB)}.
Let $\mathfrak q \subset B$ be the prime corresponding to $\xi$ and
let $\mathfrak p \subset A$ be the corresponding prime of $A$.
By assumption $\mathfrak q \in \text{Ass}_B(M \otimes_A \kappa(\mathfrak p))$,
see
Algebra, Remark \href{https://stacks.math.columbia.edu/tag/05E0}{remark bourbaki (Tag 05E0)}
and
Divisors, Lemma \href{https://stacks.math.columbia.edu/tag/02OK}{lemma associated affine open (Tag 02OK)}.
With this notation $Z = V(\mathfrak q) \subset \Spec(B)$.
In particular $f(Z) \subset V(\mathfrak p)$. Hence clearly it suffices to
prove the lemma after replacing $A$, $B$, and $M$ by
$A/\mathfrak pA$, $B/\mathfrak pB$, and $M/\mathfrak pM$.
In other words we may assume that $A$ is a domain with fraction field
$K$ and $\mathfrak q \subset B$ is an associated prime of $M \otimes_A K$.

\medskip\noindent
At this point we can use generic flatness. Namely, by
Algebra, Lemma \href{https://stacks.math.columbia.edu/tag/051T}{lemma generic flatness (Tag 051T)}
there exists a nonzero $g \in A$ such that $M_g$ is flat
as an $A_g$-module. After replacing $A$ by $A_g$ we may assume that
$M$ is flat as an $A$-module.

\medskip\noindent
In this case, by
Algebra, Lemma \href{https://stacks.math.columbia.edu/tag/05C1}{lemma post bourbaki (Tag 05C1)}
we see that $\mathfrak q$ is also an associated prime of $M$.
Hence we obtain an injective $B$-module map $B/\mathfrak q \to M$.
Let $Q$ be the cokernel so that we obtain a short exact sequence
$$
0 \to B/\mathfrak q \to M \to Q \to 0
$$
of finite $B$-modules. After applying generic flatness
Algebra, Lemma \href{https://stacks.math.columbia.edu/tag/051T}{lemma generic flatness (Tag 051T)}
once more, this time to the $B$-module $Q$, we may assume that $Q$
is a flat $A$-module. In particular we may assume the short exact
sequence above is universally injective, see
Algebra, Lemma \href{https://stacks.math.columbia.edu/tag/00HL}{lemma flat tor zero (Tag 00HL)}.
In this situation
$(B/\mathfrak q) \otimes_A \kappa(\mathfrak p')
\subset M \otimes_A \kappa(\mathfrak p')$
for any prime $\mathfrak p'$ of $A$. The lemma follows as a minimal
prime $\mathfrak q'$ of the support of
$(B/\mathfrak q) \otimes_A \kappa(\mathfrak p')$
is an associated prime of $(B/\mathfrak q) \otimes_A \kappa(\mathfrak p')$ by
Divisors, Lemma \href{https://stacks.math.columbia.edu/tag/05AH}{lemma minimal support in ass (Tag 05AH)}.
\end{proof}
\end{document}
